\cotefile{11}
\scanpage{3}
$\unsure{V|\mathcal V}$ catégorie (objets $X,Y,\ldots$) avec produits de deux \unsure{objet|objets} (\unsure{+|et} un objet final ?)

$\mathcal A\colon X\to A(X)$ foncteur cov. de $V$ dans \unsure{la|?} catégorie \deleted{\unsure{des}} \inserted{$C$} Algèbres \inserted{gr.} bivariantes, satisfaisant une condition d'échange \deleted{\unsure{entre}} \deleted{\unsure{condition}} \inserted{des $f_*$ et $g^*$} (dont le besoin exact sera \unsure{vu} ds la suite \ldots).

On va définir une nouvelle catégorie, notée $V^{\mathcal A}$, et deux foncteurs $\varphi,\psi$

\diagram{\begin{tikzcd}[ampersand replacement=\&,column sep=large]
V \arrow[r,"\varphi"] \arrow[dr,"\mathcal A"'] \& V^{\mathcal A} \arrow[r,"\psi"] \& k\text{-mod gr.}\\
\& C \arrow[ur,"\omega", "{\substack{\text{oubli}\\\text{de la structure d'\unsure{Algèbres|Algèbre|algèbre},}}}"'] \&
\end{tikzcd}}

rendant le carré écrit commutatif. De plus, le foncteur $\varphi$ induit une bijection (en fait, une identité) sur les ens d'objets, en \deleted{\unsure{d'autres termes}} \inserted{particulier}

\[
\mathrm{Ob}\,V^{\mathcal A}=\mathrm{Ob}\,V
\]

On pose.

\[
\operatorname{Hom}_{V^{\mathcal A}}(X,Y)=\mathcal A(X\times \unsure{Y|X})
\]

On définit le produit de convolution ainsi

\[
\begin{aligned}
\operatorname{Hom}_{V^{\mathcal A}}(X,Y)\times\operatorname{Hom}_{V^{\mathcal A}}(Y,Z)&\longrightarrow\operatorname{Hom}_{V^{\mathcal A}}(X,Z)\\
(u,\quad v)&\longmapsto v\circ u
\end{aligned}
\]

\[
\mathcal A(X\times Y)\times\mathcal A(Y\times Z)\longrightarrow\mathcal A(X\times Z)
\]

\diagram{\begin{tikzcd}[ampersand replacement=\&,row sep=small,column sep=small]
\&X\times Y\times Z\arrow[dl,"p_{12}"']\arrow[d,dash,"p_{13}"]\arrow[dr,"p_{23}"]\&\\
X\times Y\&X\times Z\&Y\times Z
\end{tikzcd}}

\[
\boxed{v\circ u=p_{31*}\bigl(p_{32}^{*}(\unsure{v|u})\,p_{21}^{*}(u)\bigr).}
\]

% archive sheet number: 2
\scanpage{4}
\mspage{-3.2-}

\margin{\handwritten{\unsure{11|''}}}

\[
\handwritten{\varpi}\colon\mathrm{SL}(n+1)\handwritten{\longrightarrow}\mathrm{PGL}(n)
\]

is the canonical isogeny, then with respect to the induced action of $\mathrm{SL}(n+1)$, $\mathrm{o}_{\handwritten{\mathbb P}}(1)$ admits an $\mathrm{SL}(n+1)$-linearization. This follows because $\mathrm{SL}(n+1)$ acts on the affine cone over $\handwritten{\mathbb P}_n$, compatibly with the given action of $\mathrm{PGL}(n)$ on $\handwritten{\mathbb P}_n$. This linearization induces an $\mathrm{SL}(n+1)$-linearization of $L_i$ for each $i$ (for the induced action of $\mathrm{SL}(n+1)$ ), and therefore such a linearization for each product

\[
\handwritten{\bigotimes}_{\handwritten{i=0}}^{\handwritten{m}} L_i^{\handwritten{\otimes r_i}}.
\]

On the other hand, the $(n+1)^{\mathrm{st}}$ power of the $\mathrm{SL}(n+1)$-linearization of $\underline{\mathrm{o}}_{\handwritten{\mathbb P}}(1)$, which is a linearization of $\underline{\mathrm{o}}_{\handwritten{\mathbb P}}(n+1)$, is induced, via $\handwritten{\varpi}$, from a $\mathrm{PGL}(n)$-linearization (cf. Chapter 1, \S2). Therefore, the $\mathrm{SL}(n+1)$-linearization of

\[
\handwritten{\bigotimes}_{\handwritten{i=0}}^{\handwritten{m}} L_i^{r_i}
\]

is a $\mathrm{PGL}(n)$-linearization if $n+1$ divides each $r_i$.

In any case, by means of these linearizations we can define invariant sections of all the sheaves $\handwritten{\bigotimes}_{\handwritten{i=0}}^{\handwritten{m}} L_i^{r_i}$. To construct such invariant sections, let $X_{\unsure{o|0}},\ldots,X_n$ be the canonical sections of $\underline{\mathrm{o}}_{\handwritten{\mathbb P}}(1)$ on $\handwritten{\mathbb P}_n$. Let

\[
X_i^{(j)}=p_j^*(X_i)
\]

be the induced sections of $L_j$.

\uline{Definition} \uline{3.1}: For all sequences $\handwritten{\alpha}_o,\ldots\handwritten{\alpha}_n$ of integers such that

% archive sheet number: 3
\scanpage{5}
Associativité

\diagram{\begin{tikzcd}[ampersand replacement=\&]
X\arrow[r,dash,"u"]\&Y\arrow[r,dash,"v"]\&Z\arrow[r,dash,"w"]\&T
\end{tikzcd}}

À prouver

\margin{\unsure{./}}

\[
w(vu)=(wv)u
\]

\deleted{Le premier mb est égal à \unsure{$\gamma$,}}

\diagram{\begin{tikzcd}[ampersand replacement=\&,column sep=small,row sep=large]
\&\&X\times Y\times Z\times T\arrow[dl,"\alpha'"']\arrow[dr,"\nu'"]\&\&\\
\&X\times Y\times Z\arrow[dl,"\lambda"']\arrow[d,"\mu"']\arrow[dr,"\nu","{\deleted{\unsure{X\times Z}}}"']\&\&X\times Z\times T\arrow[dl,"\alpha"'," {\deleted{\unsure{\alpha|\mu}}}"]\arrow[d,"\gamma"']\arrow[dr,"\beta"]\&\\
X\times Y\&\deleted{\unsure{X}}Y\times Z\&X\times Z\&X\times T\&Z\times T
\end{tikzcd}}

\[
\begin{aligned}
w(vu)&\overset{\text{déf\textsuperscript{n}}}{=}\gamma_*\bigl(\beta^*(w)\,\alpha^*(vu)\bigr)
\overset{\text{déf\textsuperscript{n} de }vu}{=}\gamma_*\bigl(\beta^*(w)\,\alpha^*\bigl(\nu_*(\mu^*(v)\deleted{\unsure{)}}\lambda^*(u))\bigr)\bigr)\\
&\overset{\substack{\text{car }\alpha^*\nu_*=\nu'_*\alpha'{}^*\\\boxed{\text{échange}}}}{=}
\gamma_*\Bigl(\beta^*(w)\,\nu'_*\bigl(\underbrace{\alpha'{}^*(\mu^*(v)\lambda^*(u))}_{\alpha'{}^*\mu^*(\unsure{w|v})\,\alpha'{}^*\lambda^*(u)=p_{YZ}^*(v)\,p_{X,Y}^*(u)}\bigr)\Bigr)
\end{aligned}
\]

\[
\begin{aligned}
&\underset{\text{form. de proj (att. signes !)}}{\Vert}\\
(-1)^{\overline{\unsure{w}}\,\overline{\unsure{\nu|\gamma}}_*}\,\deleted{\unsure{=}}\;\deleted{\unsure{\gamma_*|\beta_*}}\quad&
\nu'_*\bigl(\underbrace{\nu'{}^*\beta^*(w)}_{p_{ZT}^*(w)}\,p_{YZ}^*(v)\,p_{\unsure{YY|XY}}^*(u)\bigr)\\
&=p_{\unsure{X|\alpha}T*}\bigl(p_{ZT}^*(w)\,p_{YZ}^*(v)\,p_{XY}^*(u)\bigr)
\end{aligned}
\]

\diagram{\begin{tikzcd}[ampersand replacement=\&,column sep=small,row sep=large]
\&\&X\times Y\times Z\times T\arrow[dl,"n'"']\arrow[dr,"b'"]\&\&\\
\&X\times Y\times T\arrow[dl,"a"']\arrow[d,"c"']\arrow[dr,"b"]\&\&Y\times Z\times T\arrow[dl,"n"']\arrow[d,"l"]\arrow[dr,"m"]\&\\
X\times Y\&X\times T\&Y\times T\&\substack{Y\times Z\\v}\&\substack{Z\times T\\w}
\end{tikzcd}}

\[
\begin{aligned}
(wv)u&=c_*\bigl(b^*(wv)\,a^*(u)\bigr)=c_*\Bigl(\underbrace{b^*\bigl(n_*\deleted{\unsure{m|m_*}}\bigl(m^*(w)\,l^*(v)\bigr)\bigr)}_{\substack{\boxed{\text{formule d'échange}}\;\bigl\{\;n'_*\bigl(\underbrace{b'{}^*(m^*(w)\,l^*(v))}_{p_{ZT}^*(w)\,p_{YZ}^*(v)}\bigr)}}\,a^*(u)\Bigr)\\
&=c_*\Bigl(\underbrace{n'_*\bigl(p_{ZT}^*(w)\,p_{YZ}^*(v)\bigr)\,a^*(u)}_{\text{form. de pr. (att. signes !)}}\Bigr)\\
&=\deleted{c_*}\;n'_*\bigl(p_{ZT}^*(w)\,p_{YZ}^*(v)\,\underbrace{\deleted{\unsure{p}}\,\unsure{n'{}^*}a^*(u)}_{p_{XY}^*(u)}\\
&=p_{XT*}\bigl(p_{ZT}^*(w)\,p_{YZ}^*(v)\,p_{\unsure{XY}}^{*}\unsure{(}
% The remaining strokes of the argument run into the right edge.
\end{aligned}
\]

% archive sheet number: 4
\scanpage{6}
\mspage{-3.1-}

\uline{Chapter} \uline{3}. \uline{An} \uline{elementary} \uline{exemple}

The purpose of this chapter is to give, independently of the forgoing theory, an exhaustive analysis of a single special class of actions. This has two objectives: first of all, the concepts of pre-stable and stable are worked out in a representative non-trivial case. Secondly, by attacking this case directly, we can circumvent the difficulties involved in extending the previous work to semi-simple groups in characteristic $p$ (cf. Preface). And, in fact, we obtain the result over the ring of integers, hence, \uline{a} \uline{fortiori}, over any field. This will enable us later to construct the so-called arithmetic schemes of moduli of abelian varieties.

We have no use for a ground field in this chapter: all the schemes that we will consider will be of finite type over $\operatorname{Spec}(\handwritten{\mathbb Z})$. $\handwritten{\mathbb A}^{n}$, $\handwritten{\mathbb P}_n$, and $\mathrm{PGL}(n)$ stand for the usual schemes \uline{over} \uline{$\handwritten{\mathbb Z}$}. The action we shall study is the canonical action:

\[
\sigma_m:\mathrm{PGL}(n)\times(\handwritten{\mathbb P}_n)^{m+1}\handwritten{\longrightarrow}(\handwritten{\mathbb P}_n)^{m+1}
\]

\margin{\handwritten{\unsure{--}}}

i.e. of projective transformations on sequences of points in projective space.

\handwritten{\fbox}{\S1.} The first step is to obtain a large number of\handwritten{\unsure{$\mid$}}invariant affine open subsets in $(\handwritten{\mathbb P}_n)^{m+1}$. To this end, let $p_0,\ldots,p_m$ denote the projections of $(\handwritten{\mathbb P}_n)^{m+1}$ onto its $m+1$ factors. Let $L_i=p_i^*(\underline{\mathrm{o}}_{\handwritten{\mathbb P}}(1))$. Now unfortunately, $\underline{\mathrm{o}}_{\handwritten{\mathbb P}}(1)$ admits no $\mathrm{PGL}(n)$-linearization with respect to the action of $\mathrm{PGL}(n)$ on $\handwritten{\mathbb P}_n$. However, if

% archive sheet number: 5
\scanpage{7}
\uline{Unités} On définit

\[
\deleted{\unsure{\delta}}\Delta_X\in\operatorname{Hom}_{\unsure{V|\mathcal V}^{\mathcal A}\unsure{{}_{c}|{}}}(X,X)=A(X\times X)
\]

en considérant

\[
\delta_X:X\mathbin{-}X\times X\qquad(\text{diag})
\]

et définissant

\[
\Delta_X=\deleted{\unsure{\delta_{X*}A(\delta_X)}}A(\delta_X)_*(1_{A(X)})
\]

Prouvons que $\Delta_X$ est une unité bilatère pour la composition : on a

\diagram{\begin{tikzcd}[ampersand replacement=\&]
X\arrow[r,dash,"\Delta_X"]\&X\arrow[r,dash,"u"]\&Y
\end{tikzcd}}

\[
u\circ\Delta_X=p_{31*}\bigl(p_{32}^*(u)\,p_{21}^*(\Delta_X)\bigr)
\]

\diagram{\begin{tikzcd}[ampersand replacement=\&,column sep=small]
\&X\times X\times Y\arrow[dl,"p_{21}"']\arrow[d,dash]\arrow[dr,"p_{32}"]\&\\
X\times X\&X\times Y\&X\times Y
\end{tikzcd}}

\[
p_{21}^*(\Delta_X)=p_{21}^{\unsure{'*|*}}(\delta_{X*}(1))\overset{\boxed{\text{échange}}}{=}q_*(p^*(1))=q_*(1)
\]

\[
\begin{aligned}
\text{donc}\quad p_{32}^*(u)\,p_{21}^*(\Delta_X)&=p_{32}^*(\unsure{u_X|u})\,q_*(1)
\underset{\substack{\text{form. de proj.}\\\text{Attention \uline{\unsure{signe|signes}} !}}}{=}q_*\bigl(q^*(p_{32}^*(u_X))\bigr)\\
\deleted{\unsure{\text{or}}}\quad&=q_*(u)\qquad\text{car }p_{32}q=\mathrm{id}_{X\times Y}
\end{aligned}
\]

\diagram{\begin{tikzcd}[ampersand replacement=\&]
X\times Y\arrow[r,"\delta_X\times\mathrm{id}_Y=q"]\arrow[d,"{\unsure{p|p_*}}"']\&X\times X\times Y\arrow[d,"p_{21}"]\\
X\arrow[r,"\delta_X"']\&X\times X
\end{tikzcd}}

\margin{\fbox{\parbox{2cm}{Signes $(-1)^{ux}$ ; pas de signes pour \unsure{constater|calculs} que $\Delta_X\circ v=v$ $v:Y\mathbin{-}X$}}}

\[
\begin{aligned}
\text{donc}\quad u\circ\Delta_X&=p_{31*}\,q_*(u)=u\\
&\text{car }\deleted{q}\,p_{31}q=\mathrm{id}
\end{aligned}
\]

On prouve de façon analogue que $\Delta_X\circ u=u$.

\uline{Donc on trouve bien une catégorie $V^{\mathcal A}$}

[ \uline{NB} Pour définir la composition on n'a eu besoin que des $p_*,p^*$ pour projections $X\times Y\underset{\deleted{\unsure{\leftarrow|\rightarrow}}}{\longrightarrow}X$ ; mais pour \deleted{définir} construire les unités on a utilisé aussi $q_*,q^*$ pour les imm. diagonales et leurs produits par un $Y$ \ldots{} ]

% archive sheet number: 6
\scanpage{8}
\mspage{-2.33-}

then:

\[
\begin{aligned}
\handwritten{\unsure{\nu|v}}^{M\handwritten{\otimes} f^*L^n}(x,\handwritten{\delta})&=\handwritten{\nu}^{M\handwritten{\otimes}f^*L^{\handwritten{n_0}}}(x,\handwritten{\delta})+\handwritten{\nu}^{f^*L^{\handwritten{n-n_0}}}(x,\handwritten{\delta})\\
&=\handwritten{\nu}^{M\handwritten{\otimes}f^*L^{\handwritten{n_0}}}(x,\handwritten{\delta})+(n-n_0)\handwritten{\nu}^{L}(f(x),\handwritten{\delta}).
\end{aligned}
\]

First of all, apply the first part of Proposition 2.17 to $\handwritten{\nu}^{M\handwritten{\otimes}f^*L^{\handwritten{n_0}}}$. We conclude that, for some $K$:

\[
\handwritten{\nu}^{M\handwritten{\otimes}f^*L^{\handwritten{n_0}}}(x,\handwritten{\delta})\geq-K.
\]

Secondly, apply the second part of Proposition 2.17 to $\handwritten{\nu}^L(y,\handwritten{\delta})$, for $y\in S=Y^s_{(\unsure{o|0})}(L)_k$. By\unsure{$\shortparallel$|n} Theorem 2.1, all these numbers are strictly positive. Therefore, by the Proposition:

\[
\handwritten{\nu}^L(y,\handwritten{\delta})\geq\epsilon>0
\]

for all $y\in Y^s_{(o)}(L)_k$, and $\handwritten{\delta}\in\Delta(G)$. Therefore, if $n\geq\handwritten{\left[}n_0+\frac{K}{\epsilon}+1\handwritten{\right]}$, we conclude

\[
\handwritten{\nu}^{M\handwritten{\otimes}f^*L^n}(x,\handwritten{\delta})>0
\]

whenever $x\in f^{-1}\handwritten{\{}Y^s_{(o)}(L)\handwritten{\}}_k$. Therefore, by Theorem 2.1, all such $x$ are in $X^{s\unsure{{}!|{}'|{}}}_{(o)}(M\handwritten{\otimes} f^*L^n)$. \uline{QED}

% archive sheet number: 7
\scanpage{9}
De \unsure{$\bar{\mathrm m}$|même|fait}, pour la \deleted{\unsure{définition|définit}} vérification de l'associativité, on n'a utilisé la formule de projection et la formule d'échange que ds le cas

\diagram{\begin{tikzcd}[ampersand replacement=\&,row sep=small]
\& X\times Y\times \unsure{P|p}\arrow[dl]\arrow[dr]\&\\
X\times P\arrow[dr,dash]\&\&\deleted{X}Y\times P\arrow[dl,dash]\\
\&P\&
\end{tikzcd}}

, mais pour vérifier que les $\unsure{\delta|\Delta}_X$ sont des unités, on \deleted{\unsure{en}} avait besoin aussi dans un cas

\diagram{\begin{tikzcd}[ampersand replacement=\&,row sep=small]
\& X\times Y\arrow[dl]\arrow[dr,hook]\&\\
X\arrow[dr,hook,"\text{diag.}"']\&\&X\times X\times Y\arrow[dl,"\text{proj}"']\\
\&X\times X\&
\end{tikzcd}}

\noindent\rule{2.7cm}{0.3pt}\par\nobreak\vspace{-0.8\baselineskip}\noindent\rule{3cm}{0.3pt}

Définissons le foncteur $\varphi$, \unsure{donc|en} définissons pour tt $X,\unsure{Z|Y}\in\mathrm{Ob}\,\unsure{V|\mathcal V}$ une application

\[
\varphi_{X,Z}:\operatorname{Hom}_{V}(X,Y)\longrightarrow\operatorname{Hom}_{V^{\mathcal A}}(X,Y)=A(X\times Y)
\]

Soit donc $f\unsure{{:}|\in}X\longrightarrow Y$ dans $V$, on en déduit un hom graphe

\[
\unsure{\Gamma|\mathscr T}_f=(\mathrm{id}_X,f)\unsure{. |{:}}X\unsure{\mathbin{-}|\longrightarrow}X\times Y
\]

\[
\text{d'où}\qquad A(\Gamma_f)_*:A(X)\mathbin{-}A(X\times Y)
\]

Posons alors

\[
\deleted{\unsure{A}}\ \varphi_{X,Z}(f)=A(\Gamma_f)_*(1)
\]

% archive sheet number: 8
\scanpage{10}
\mspage{-2.32-}

$G$ on $X_1$ extending the given one on $X$, and there is an $M_1\in\operatorname{Pic}^{G}(X_1)$, ample on $X_1$, which \inserted{\handwritten{extends}} $M^N\handwritten{\otimes}f^*L^{N_0N}$ for some $N$. (This $N$ may be taken as the same $N$ as above). Then we have the diagram:

\diagram{\begin{tikzcd}[ampersand replacement=\&]
X\handwritten{\arrow}[r,"{(I_X,I_Y\circ f)}"]\handwritten{\arrow}[d,"f"']\&X_1\times Y_1\handwritten{\arrow}[d,"{{}_{\unsure{-|{}}}\unsure{p|P}_2}"]\\
Y\handwritten{\arrow}[r,hook,"I_Y"']\&Y_1
\end{tikzcd}}

Now if the result is proven when $X$ and $Y$ are proper, we can apply it to $p_2:X_1\times Y_1\handwritten{\longrightarrow} Y_1$, and the sheaves $M_1\handwritten{\otimes}\underline{\mathrm{o}}_{Y_1}$ on $X_1\times Y_1$ and $L_1$ on $Y_1$. Then

\[
\begin{aligned}
f^{-1}\handwritten{\{}Y^s_{(\unsure{o|0})}(L)\handwritten{\}}&=f^{-1}\handwritten{\{}I_Y^{-1}\handwritten{\{}Y^s_{1,(o)}(L_1)\handwritten{\}}\handwritten{\}}\\
&=(I_X\unsure{{;}|{}_{j}|{,}}I_Y\circ f)^{-1}\handwritten{\{}p_2^{-1}\handwritten{\{}Y^s_{1,(o)}(L_1)\handwritten{\}}\handwritten{\}}\\
&\handwritten{\subset}(I_X\mathpunct{;}I_Y\circ f)^{-1}\handwritten{\{}(X_1\times Y_1)^s_{(o)}(M_1\handwritten{\otimes}L_1^n)\handwritten{\}}\ ;\ \text{if }n\geq n_0\\
&\handwritten{\subset}X^s_{(o)}\handwritten{\bigl[}(I_X,I_Y\circ f)^*(M_1\handwritten{\otimes}L_1^n)\handwritten{\bigr]}\qquad\text{by Prop. 1}\quad,\\
&=X^s_{(o)}\handwritten{\bigl[}M^N\handwritten{\otimes} f^*L^{NN_0+Nn}\handwritten{\bigr]}\\
&=X^s_{(o)}\handwritten{\bigl[}M\handwritten{\otimes} f^*L^{N_0+n}\handwritten{\bigr]}.
\end{aligned}
\]

Now, if $X$ and $Y$ are proper over $k$, we note again that $M\handwritten{\otimes} f^*L^n$ is ample on $X$ for $n\geq n_0$. Therefore, we can apply Theorem 2.1 to the computation of $X^s_{(o)}(M\handwritten{\otimes} f^*L^n)$. But if $x\in X_k$, $\handwritten{\delta}\in\Delta(G)$,

% archive sheet number: 9
\scanpage{11}
Il faut prouver qu'on a bien un foncteur.

a) $\varphi_{X,X}(\mathrm{id}_X)=\Delta_X$. En effet, \deleted{\unsure{le produit est}} \inserted{c'est la \uline{définition}} de $\Delta_X$ !

\diagram{\begin{tikzcd}[ampersand replacement=\&]
\text{b)}\quad X\arrow[r,dash,"f"]\&Y\arrow[r,dash,"g"]\&Z\qquad\varphi(g)\varphi(f)=\varphi(gf)
\end{tikzcd}}

Le premier mb est

\diagram{\begin{tikzcd}[ampersand replacement=\&,column sep=small,row sep=large]
\&X\times Z\arrow[dl,"\mu'"']\arrow[dr,"\lambda'"]\&\&X\times Y\arrow[dl,"{\unsure{\lambda^{?}|\lambda}=\mathrm{id}_X\times\unsure{\Gamma|\mathscr T}_g}"']\arrow[dr,"\mu"]\&\\
X\arrow[dr,"\Gamma_f"']\&\&X\times Y\times Z\arrow[dl,"p_{21}"']\arrow[d,"p_{31}"]\arrow[dr,"p_{32}"]\&\&Y\arrow[dl,"\Gamma_g"]\\
\&X\times Y\&X\times Z\&Y\times Z\&
\end{tikzcd}}

\[
p_{31*}\Bigl(
\underbrace{p_{32}^*(\varphi(g))}_{\substack{\Vert\\p_{32}^*(\Gamma_{g*}(1))\\\Vert\\\deleted{\unsure{\mu_*|p_*}}\lambda_*(\mu^*(1))=\lambda_*(1)}}
\underbrace{p_{21}^*(\varphi(f))}_{\substack{\Vert\\p_{21}^*\Gamma_{f*}(1)\\\Vert\\\lambda'_*(\mu'{}^*(1))=\lambda'_*(1)}}
\]

\[
\varphi(g)\varphi(f)=p_{31*}\bigl(\lambda^{\unsure{+|*|{}}}_*(1)\,\lambda'_*(1)\bigr)\qquad\text{Mais}
\]

\diagram{\begin{tikzcd}[ampersand replacement=\&,column sep=small,row sep=large]
\&\&X\arrow[dl,hook,"\Gamma_{gf}"']\arrow[dr,hook,"\Gamma_f"]\&\&\\
\&X\times Z\arrow[dl,"\mu'=p_X"']\arrow[dr,hook,"{\lambda'=\deleted{\unsure{I}}\Gamma_f\times Z}"']\&\text{(cartésien)}\&X\times Y\arrow[dl,hook,"{\lambda=\mathrm{id}_X\times\Gamma_g}"]\arrow[dr,"\mu=p_Y"]\&\\
X\arrow[dr,hook,"\Gamma_f"']\&\&X\times Y\times Z\arrow[dl,"p_{21}"']\arrow[d,"p_{31}"]\arrow[dr,dash,"p_{32}"]\&\&Y\arrow[dl,hook,"\Gamma_g"]\\
\&X\times Y\&X\times Z\&Y\times Z\&
\end{tikzcd}}

\[
\begin{aligned}
\lambda_*(1)\,\unsure{{}_*|{}}\lambda'_*(1)&=\lambda_*\bigl(\underbrace{\lambda^*(\lambda'_*(1))}_{\substack{\Vert\\\Gamma_{f*}(\Gamma_{gf}^*(1))=\Gamma_{f*}(1)}}\bigr)\quad\boxed{\text{form. proj}}\\
&=\lambda_*\Gamma_{f*}(1)=\unsure{\nu|v}_*(\unsure{1|1/1})
\end{aligned}
\]

où $\nu:X\mathbin{-}X\times Y\times Z$ est donné par $\nu(x)=(x,f(x),gf(x))$

\[
\begin{aligned}
x&\mapsto (x,gf(x))\unsure{\vdash|\mapsto|\leftarrow} (x,f(x),gf(x))\\
x&\vdash (x,f(x))\vdash (x,f(x),\unsure{gf(x))}
\end{aligned}
\]

\[
\begin{aligned}
&\deleted{\unsure{x\mapsto{}}}\ (x,f(x),z)\\
&(x,y,g(y))\\
&(x,f(x),gf(x))
\end{aligned}
\]

D'où

\[
\varphi(g)\varphi(f)=p_{31*}\bigl(\nu_*(1)\bigr)=\Gamma_{gf*}(1)=\varphi(gf)\qquad\text{OK.}
\]

\noindent\rule{3.5cm}{0.3pt}\par\nobreak\vspace{-0.8\baselineskip}\noindent\rule{3.6cm}{0.3pt}

Définissons le foncteur $\psi$, \inserted{avec $\psi(X)=A(\unsure{X|x})$,} i.e. pour $u\in\operatorname{Hom}_{\unsure{V|\mathcal V}^{\mathcal A}}(X,Y)=A(X\times Y)$, et $x\in A(X)$, définissons $u(x)\in A(Y)$ par \unsure{produit|la formule} \deleted{\unsure{$u(x)$}}

\[
u(x)=p_{2*}(u\,p_1^*(x))
\]

Prouvons que l'on a bien une \uline{loi fonctorielle}.

% archive sheet number: 10
\scanpage{12}
\mspage{-2.31-}

minimized is invariant under the transformation $z\handwritten{\longrightarrow\alpha}z$, $\handwritten{\alpha}>0$, and the set of rays in $\handwritten{\unsure{\Gamma|\Pi|\mathscr T}^{\mathbb R}}(T)$ is \uline{compact}. We must check that this ``inf'' is actually attained by a rational point $z$. But each linear functional $\handwritten{\ell}_i$ is rational, and the inner product $<x,y>$ on $\handwritten{\Gamma^{\mathbb R}}(T)$ is rational: then it is an exercise in elementary calculus to check that this is so. \uline{QED}

An application of this is:

\uline{Proposition} \uline{2.18}: Let $G$ be a reductive group acting on algebraic schemes $X$ and $Y$. Let $f:X\handwritten{\to}Y$ be a $G$-linear morphism, let $L\in\operatorname{Pic}^G(Y)$ and $M\in\operatorname{Pic}^G(X)$. Then if $M$ is relatively ample for $f$ and $L$ is ample on $Y$, there is an $n_0$ such that:

\[
n\geq n_0\handwritten{\Longrightarrow}X^s_{(\unsure{o|0})}(M\handwritten{\otimes}f^*L^n)\handwritten{\supset} f^{-1}\handwritten{\{}Y^s_{(o)}(L)\handwritten{\}}.
\]

\uline{Proof}: We shall first reduce this Proposition to the case where $X$ and $Y$ are proper over $k$. First of all, there is an open immersion $I_Y:Y\handwritten{\subset}Y_1$, where $Y_1$ is proper over $k$, and there is an action of $G$ on $Y_1$ extending the given one of $G$ on $Y$, and there is an $L_1\in\operatorname{Pic}^G(Y_1)$, ample on $Y_1$ which extends $L^N$ for some $N$, such that:

\[
Y^s_{(o)}(L)=I_Y^{-1}\handwritten{\{}Y^s_{1,(o)}(L_1)\handwritten{\}}
\]

(cf. Proposition 1.\handwritten{8} ). Secondly, for some $N_0$, $M\handwritten{\otimes}f^*L^{N_0}$ is ample on $X$ by EGA2, \handwritten{4.6}:13, (ii). Then there is also an open immersion $I_X:X\handwritten{\subset}X_1$, where $X_1$ is proper over $k$, and there is an action of

% archive sheet number: 11
\scanpage{13}
a) Si $X=Y$, prouvons que $\Delta_X(x)=x$. On a

\[
\begin{aligned}
\Delta_X(x)&=p_{2*}(\Delta_X\,p_1^*(x))
=p_{2*}\Bigl(\underbrace{\delta_*(1)\,p_1^*(x)}_{\substack{\Vert\quad\text{form. de proj}\\\delta_*(\underbrace{1.\delta^*p_1^*(x)}_{x})}}\Bigr)\\
&\unsure{=|{}}\ =p_{2*}\,\delta_*(x)=x\qquad\text{OK.}
\end{aligned}
\]

\diagram{\begin{tikzcd}[ampersand replacement=\&,row sep=small]
\&X\arrow[d,"\delta"']\&\\
\&X\times X\arrow[dl,"p_1"']\arrow[dr,dash,"p_2"]\&\\
X\&\&X
\end{tikzcd}}

b) Prouvons $v(u(x))=(v\circ u)(x)$.

\diagram{\begin{tikzcd}[ampersand replacement=\&]
X\arrow[r,"u"]\&Y\arrow[r,"v"]\&Z
\end{tikzcd}}

\diagram{\begin{tikzcd}[ampersand replacement=\&,column sep=small]
\&\&X\times Y\times Z\arrow[dl,"{\lambda=p_{XY}}"']\arrow[dr,"{\mu=p_{YZ}}"]\&\&\\
\&X\times Y\arrow[dl,"p_1"']\arrow[dr,"p_2"]\&\&Y\times Z\arrow[dl,"q_1"']\arrow[dr,"q_2"]\&\\
\deleted{\unsure{Y}}X\&\&Y\&\&Z
\end{tikzcd}}

\[
\begin{aligned}
v(u(x))&=\deleted{\unsure{q_1^*(v|p_1^*}}\,q_{2*}\bigl(v\,q_1^*u(x)\bigr)\\
&=q_{2*}\Bigl(v\,\underbrace{q_1^*\bigl(p_{2*}(u\,p_1^*(x))\bigr)}_{\substack{\Vert\quad\boxed{\text{échange}}\\\mu_*(\lambda^*(u\,p_1^*(x)))\\\Vert\\\mu_*(\lambda^*(u)\,p_X^*(x))}}\Bigr)\\
&\deleted{\unsure{=}}\\
&\underset{\text{proj, attention signes}}{=}q_{\unsure{2|n}*}\bigl(\mu_*(\mu^*(v)\,\lambda^*(u)\,p_X^*(x)\bigr)\\
&=p_{Z*}\bigl(p_{YZ}^*(v)\,p_{XY}^*(u)\,p_{\unsure{X|\alpha}}^*(x)\bigr)
\end{aligned}
\]

On a d'autre part

\diagram{\begin{tikzcd}[ampersand replacement=\&,column sep=small,row sep=small]
\&X\times Y\times Z\arrow[dl,"p_{XY}"']\arrow[d,"p_{YZ}"']\arrow[dr,"p_{XZ}"]\&\\
X\times Y\&Y\times Z\&X\times Z\arrow[dll,"r_1"']\arrow[d,"r_2"]\\
X\&\&Z
\end{tikzcd}}

\[
\begin{aligned}
(v\circ u)(x)&=r_{2*}\bigl((v\circ u)\,r_1^*(x)\bigr)
=r_{2*}\Bigl(\underbrace{\cdot p_{XZ*}\bigl(p_{YZ}^*(v)\,p_{XY}^*(u)\bigr)\,r_1^*(x)}_{p_{XZ*}\bigl(p_{YZ}^*(v)\,p_{XY}^*(u)\,\underbrace{p_{XZ}^*r_1^*(x)}_{p_X^*(x)}\bigr)}\Bigr)\\
&=p_{Z*}\bigl(p_{YZ}^*(v)\,p_{XY}^*(u)\,p_X^*(x)\bigr)
\end{aligned}
\]

OK.

% archive sheet number: 12
\scanpage{14}
\mspage{-2.30-}

\uline{Proof}: Let $T\handwritten{\subset}G$ be a maximal torus. Let $<x,y>$ be the form on $\handwritten{\unsure{\Gamma|\Pi|\mathscr T}^{\mathbb R}}(T)$ which defines the norm $\Vert\lambda\Vert$. Let $\handwritten{\ell}_1,\ldots,\handwritten{\ell}_n$ be the linear functionals on $\handwritten{\Gamma^{\mathbb R}}(T)$ given by Proposition 2.14. For any $x\in X_k$, let $I(x)$ be the subset $\handwritten{\{}1,\ldots,n\handwritten{\}}$ given by Proposition 2.14. To estimate $\handwritten{\unsure{\nu|v}}^L(x,\handwritten{\delta})$, say $\handwritten{\delta}=\Delta(\lambda)$. There is always some $\handwritten{\gamma}\in G_k$ such that $\handwritten{\gamma}\circ\lambda\circ\handwritten{\gamma}^{-1}$ is a 1-PS of $T$. Then:

\[
\begin{aligned}
\handwritten{\nu}^L(x,\handwritten{\delta})&=\frac{\handwritten{\mu}^{L}(x,\lambda)}{\Vert\lambda\Vert}\\
&=\frac{\handwritten{\mu}^{L}(\sigma(\handwritten{\gamma},x),\handwritten{\gamma}\circ\lambda\circ\handwritten{\gamma}^{-1})}{\Vert\handwritten{\gamma}\circ\lambda\circ\handwritten{\gamma}^{-1}\Vert}\\
&=\max\handwritten{\left\{}\frac{\handwritten{\ell}_i(\handwritten{\overline}{\handwritten{\gamma}\circ\lambda\circ\handwritten{\gamma}^{-1}})}{\sqrt{\langle\handwritten{\overline}{\handwritten{\gamma}\circ\lambda\circ\handwritten{\gamma}^{-1}},\handwritten{\overline}{\handwritten{\gamma}\circ\lambda\circ\handwritten{\gamma}^{-1}}\rangle}}\ \handwritten{\middle|}\ i\in I(\sigma(\handwritten{\gamma},x))\handwritten{\right\}}
\end{aligned}
\]

But, for each $i$, $\handwritten{\left|}\handwritten{\ell}_i(z)\handwritten{\right|}\big/\sqrt{<z,z>}$ is bounded when $z$ varies over all non-zero points of $\handwritten{\Gamma^{\mathbb R}}(T)$. Therefore $\handwritten{\nu}^L(x,\handwritten{\delta})$ is bounded when $x$ and $\handwritten{\delta}$ vary arbitrarily. Now suppose for $x\in S\handwritten{\subset}X_k$, we seek to minimize $\handwritten{\nu}^L(x,\handwritten{\delta})$. According to the above,

\[
\inf_{\substack{\handwritten{\delta}\in\Delta(G)\\x\in S}}\handwritten{\nu}^{L}(x,\handwritten{\delta})
=\inf_{\substack{\handwritten{\gamma}\in G_k\\x\in S}}
\handwritten{\left[}\inf_{z\in\handwritten{\Gamma^{\mathbb Q}}(T)}
\handwritten{\left[}\max_{i\in I(\sigma(\handwritten{\gamma},x))}
\frac{\handwritten{\ell}_i(z)}{\sqrt{<z,z>}}\handwritten{\right]}\handwritten{\right]}
\]

But the first ``inf'' is essentially a minimum over the finite set of subsets $I\handwritten{\subset\{}1,\ldots,n\handwritten{\}}$ which occur as subsets $I(\sigma(\handwritten{\gamma},x))$. Therefore, this ``inf'' will always be attained for some $x$ and $\handwritten{\gamma}$. As for the inner ``inf'', if we replace the variable \inserted{\handwritten{by}} $z\in\handwritten{\Gamma^{\mathbb R}}(T)$, then the ``inf'' will certainly be attained for some $z$ - for the expression to be

% archive sheet number: 13
\scanpage{15}
Commutativité du carré de foncteurs. Il suffit de le voir sur les flèches, donc : prouver pour $f:X\to Y$ dans $\unsure{V|\mathcal V}$, et $x\in A(\unsure{\alpha|X|x})$

\[
f_*(x)=\varphi(f)(x)
\]

Or on a

\diagram{\begin{tikzcd}[ampersand replacement=\&,row sep=small]
\&X\arrow[d,"{\unsure{\Gamma|\mathscr T}_f}"]\&\\
\&X\times Y\arrow[dl,"p_1"']\arrow[dr,dash,"p_2"]\&\\
X\&\&Y
\end{tikzcd}}

\[
\begin{aligned}
\varphi(f)(x)&=p_{2*}\bigl(\varphi(f)\,p_1^*(x)\bigr)
=p_{2*}\Bigl(\underbrace{\Gamma_{f*}(1)\,p_1^*(x)}_{\Gamma_{f*}(\underbrace{\Gamma_f^*p_1^*(x)}_{x})}\Bigr)\\
&=p_{2*}\,\Gamma_{f*}(x)=x\qquad\text{OK.}
\end{aligned}
\]

\uline{\deleted{Question}} \inserted{Remarque} Comment interpréter le foncteur contravariant $V^{\unsure{o|2|\circ}}\mathbin{-}k\text{-mod gr}$ défini par $\mathcal A$, en composant $\mathcal A$ et le foncteur $\omega^*:C^o\longrightarrow k\text{-mod gr}$ (associant à $(f_*,\unsure{f^*|f})$ l'application $f^*$) ? En d'autres termes, comment interpréter $A(f)^*:A(Y)\mathbin{-}A(X)$, si $f:Y\longrightarrow X$ ? On \unsure{devine|démontre} que c'est défini par $\varphi(f)^*\in\operatorname{Hom}_{V^{\mathcal A}}(Y,X)=A(Y\times X)$, \unsure{déduis|déduit} de $\varphi(f)\in\unsure{H|A}(X\times Y)$ par l'isom can. qu'on devine. Vérification comme dessus.

On définit un \unsure{contr. foncteur|contre-foncteur|anti-foncteur|iso. foncteur} canonique

\[
\unsure{s|S}:V^{A^{\unsure{o|\circ}}}\mathbin{-}V^{\mathcal A^{o}},
\]

% slanted
\margin{Vérifier que c'est \uline{multiplicatif} !}

induisant l'identité sur les objets, par

\[
u\mapsto{}^{\unsure{s|t}}u:\operatorname{Hom}_{V^{\mathcal A}}(X,Y)=A(X\times Y)\mathbin{-}\operatorname{Hom}_{V^{\mathcal A}}(Y,X)=A(Y\times X)
\]

% archive sheet number: 14
